\scanpage{67}
Pour les foncteurs en sens inverse, allant de $C_{1S}=\operatorname{For}(\mathrm{GP}(1)_S)$ vers les autres catégories, on peut, soit passer par le torseur $R=\underline{\operatorname{Isom}}(G_{0S},G)$ puis tordre $\mathcal X_{0S}$
\[
G\unsure{\longmapsto|\vdash}\underline{\operatorname{Isom}}(G_{0S},G)\wedge_{G_{0S}}\mathcal X_{0S}\ ,\tag*{(21)}
\]
soit trouver, dans divers cas, des descriptions plus directes. Par exemple
\[
\begin{aligned}
G&\longmapsto\widetilde G=\text{Rev.\ simplement connexe}\\
 &\hspace{6em}\text{de $G$, défini comme}\\
 &\hspace{6em}\text{le schém. en groupes \unsure{ss|SS}}\\
 &\hspace{6em}\text{lisse à fibres \deleted{\unsure{géométr.|géom. connexes}} simplement}\\
 &\hspace{6em}\text{connexes, muni d'}\\
 &\hspace{6em}\text{une isog. $\widetilde G\to G$}\\
 &\hspace{5em}\text{-- lequel est défini de}\\
 &\hspace{6em}\text{façon \uline{unique} \ldots.}
\end{aligned}\tag*{(22)}
\]

Le foncteur vers les algèbres de Lie est
\[
\begin{aligned}
G&\longmapsto\mathfrak g=\underline{\operatorname{Lie}}(G)\\
\operatorname{For}(\mathrm{GP}(1)_S)&\xrightarrow{\unsure{\sim}}\operatorname{For}(\mathfrak{gp}(1)_S)
\end{aligned}\tag*{(23)}
\]
\[
\begin{aligned}
G&\longrightarrow\widetilde{\mathfrak g}=\underline{\operatorname{Lie}}(\widetilde G)\\
\operatorname{For}(\mathrm{GP}(1)_S)&\xrightarrow{\unsure{\sim}}\operatorname{For}(\mathfrak{sl}(2)_S)
\end{aligned}\tag*{(24)}
\]

Le foncteur
\[
G\longmapsto(E,q,\omega)
\]
peut se factoriser par $G\to\widetilde{\mathfrak g}$, cf plus bas.
% archive sheet number: 66
